\documentclass[11pt]{article}

\usepackage[margin=1in]{geometry}
\usepackage{booktabs}
\usepackage{array}
\usepackage{xurl}
\usepackage{hyperref}
\usepackage{natbib}
\usepackage{amsmath}
\usepackage{amssymb}

\newcommand{\code}[1]{\texttt{#1}}

\title{Artificial Economies for Learning Agents:\\
Testing Whether Institutions Stay Robust Across Learning Architectures}
\author{
  Thomas Prada\\
  Universidad Nacional de Colombia\\
  \texttt{tpradae@unal.edu.co}
}
\date{Draft v0.3 -- \today}

\begin{document}
\maketitle

\begin{abstract}
This paper introduces a code-first experimental platform for studying whether
economic mechanisms remain robust when the strategic agents inside them are
adaptive learners rather than equilibrium solvers. The platform implements a
shared \code{World}/\code{Agent}/\code{Institution} interface, repeated
multiseed evaluation, exploitability tests, and interchangeable mind classes
ranging from random and tabular Q-learning agents to PyTorch DQN, PPO,
decorrelated independent-DQN, and centralized-critic agents. The first
validated environment is a repeated two-firm pricing arena with six regulatory
mechanisms. In this world, price caps robustly reduce exploitability across
mind classes, but their effect on collusion is metric- and
architecture-dependent. A new paired-seed audit shows that DQN's positive
profit-under-cap result is not visible in 1000-step traces; it emerges with
training budget, becomes mostly positive by 10,000 steps, and is positive in
the 40,000-step full table. The second environment, Resource Island, extends
the same interface to a spatial inventory economy and reveals a different
lesson: institution conclusions are only meaningful when the world actually
exercises the institutional channel. Early Resource Island versions
mechanically under-tested property rights and trade price controls; a hardened
v1 now creates contested resources, unequal trades, and specialization
pressure. Three additional environments extend the suite: Auction House
provides a benchmarked auction-design testbed, Public Goods exposes common-pool
extraction and contribution incentives, and Labor Market tests learned
reporting under deferred acceptance. The current contribution is both empirical
and methodological: it reports initial evidence that institutional robustness
changes across learning architectures, and it documents the validation
discipline needed to distinguish real economic findings from inactive-mechanism
or short-horizon artifacts.
\end{abstract}

\section{Introduction}

Economic mechanisms are usually evaluated against a model of strategic behavior:
equilibrium prices, truthful bidding, efficient allocation, or some other
theoretical benchmark. That evaluation is only as good as the behavioral model
behind it. This paper sets out to find out what happens when institutions designed around classical economic assumptions are placed in front of learning agents; which guarantees survive, which break, and does the break depend on the
learning architecture?

We answer this with five small, controlled economic worlds, a pricing duopoly (Pricing Arena), a
scarce-resource trade economy (Resource Island), a sealed-bid auction (Auction House), a common-pool resource
game (Public Goods), and a two-sided labor match (Labor Market), each paired with a classical benchmark and a
shared ladder of learning architectures: random behavior, tabular Q-learning,
deep Q-networks, proximal policy optimization, independent multi-agent
Q-learning, and centralized-critic training. These architectures are not a scale
from weaker to stronger. The question at stake in each world is not which
architecture performs best, but which of the institution's guarantees hold up as
the behavioral assumptions behind the benchmark are changed.

Across the five worlds, the same failure modes recur in different forms: an
institution can look effective on one metric while a different strategic channel
stays open; an apparent effect can depend on training horizon and reverse sign with more training budget; a learning architecture can fail to recover a
benchmark that classical theory says is optimal; and an institution can raise a
headline welfare number while quietly eroding the property, stability,
sustainability, truthfulness, that the institution actually exists to protect. The next section sets out, world by
world, the classical benchmark each environment owes and what claim boundary
follows from it, while sections 4 through 8 report where each of these failure modes appears, and where
the evidence is instead null or still provisional. 

\section{Related Work}

This paper connects five economic literatures to a shared reinforcement-learning testbed. Each world below is built to test one specific classical prediction. The Pricing Arena is motivated by work on algorithmic collusion: repeated-game Q-learning agents can learn to sustain prices above the competitive level, without any explicit communication between them \citep{calvano2020,calvano2021}. The exploitability check used in this paper follows the same concern one step further. A policy's average performance across a run does not show whether that policy is robust to a deviating opponent, so this paper measures both. Auction House is based on classical auction theory. In a second-price auction, bidding one's true value is a dominant strategy, meaning it is optimal regardless of what other bidders do \citep{vickrey1961}. In a first-price auction, by contrast, the optimal strategy is to bid below one's true value (shading), and reserve prices create a known tradeoff between the seller's expected revenue and the efficiency of the final allocation \citep{myerson1981,milgrom2004}.

Public Goods and Resource Island are both based on common-pool resource and public-goods theory. The standard prediction in this literature is that, absent some form of monitoring, reputation, or locally appropriate governance, a shared resource will be overused or a public good will be underfunded, because each agent's private incentive does not account for the cost imposed on everyone else \citep{samuelson1954,ostrom1990}. Labor Market is based on two-sided matching theory. Under the worker-proposing deferred acceptance algorithm, the resulting match is guaranteed to be stable (no worker and employer who are not matched to each other would both prefer to be), and truthful reporting of preferences is guaranteed to be optimal for the proposing side \citep{galeshapley1962,rothsotomayor1990,roth1984}. These are proven properties of the mechanism, not empirical tendencies, which is what makes them useful as benchmarks.

The learner suite is built from five standard reinforcement learning methods: tabular Q-learning, deep Q-networks (DQN) \citep{sutton2018,mnih2015}, proximal policy optimization (PPO) \citep{schulman2017}, independent Q-learning \citep{tan1993}, and centralized-critic training \citep{lowe2017}. What is different here is that the suite is used to test assumptions, not to compare performance. For example, second-price truthfulness and deferred-acceptance strategy-proofness are proofs about a single, fixed learner responding to a fixed mechanism; independent Q-learning and centralized-critic training instead put multiple agents in motion at the same time, which is not the setting those proofs were written for.

\section{Theory Obligations}

Every world in this paper owes something to the theory it is built on: a specific prediction it must be checked against, a specific set of measurements that follow from that prediction, and a specific boundary on what the resulting evidence can support. This section states that obligation for each world before any learned-agent result is discussed. The obligation exists because passing tests and producing a complete set of output files is not the same thing as having tested the institution a world claims to test.

The five worlds do not owe the same kind of proof. Grouping them by the strength of their classical prediction gives three tiers: proven mechanism properties, where a learned agent's behavior can be checked directly against a known correct answer; bracketed or reference-point theory, where exact values exist but no single predicted outcome for the object under test; and activation, where no external benchmark exists at all and the obligation is instead to show the institution was ever placed under real economic pressure.

Auction House and Labor Market belong to the first tier. In Auction House, the classical predictions are exact: truthful bidding in a second-price auction, bid shading in a first-price auction, and a known revenue-efficiency tradeoff under a reserve price. Because the correct behavior is known, this world measures distance from it directly, as revenue, allocative efficiency, bidder surplus, ex-post regret, over- and underbidding, shading distance, and the shape of the learned bid curve against the truthful or shaded reference, rather than raw payoff alone. The resulting claim is bounded to benchmark deviation under different learners; convergence to auction equilibrium is not asserted. In Labor Market, a lower truthful-report rate under some learner is not by itself evidence of profitable manipulation, given the proposing side's strategy-proofness guarantee. Match rate, blocking-pair stability, truthful-report rate, and welfare are therefore reported together, since a learner can trade one against another, and any claim about manipulation is bounded to the worker-proposing, canonical version of the mechanism implemented here.

Pricing Arena and Public Goods belong to the second tier. Pricing Arena has two exact static prices, the Nash price and the joint-profit price, but no closed-form prediction for behavior in the repeated game itself, only a range bounded by the two. A collusion claim in this world is therefore only as credible as the decomposition behind it: price, profit, quantity, welfare, and exploitability are reported together and checked against each other, and the resulting claim is limited to finite-run, regulatory evidence in a stylized duopoly, not a theorem or a real-market antitrust result. Public Goods is bracketed rather than pointed: theory predicts underprovision when private incentives diverge from the social good, with free-riding as a lower bound and the social optimum as an upper one, and nothing more specific in between. The obligation this creates is to keep the welfare number honest against the underlying resource state, since an institution can raise measured welfare through reward accounting without changing the incentive to extract or contribute at all. Contribution, extraction, sustainability, inequality, and tax revenue are therefore reported as separate quantities from welfare, so that an accounting effect is not mistaken for a real one.

Resource Island belongs to the third tier alone. It carries no external closed-form benchmark, so its obligation is not how close a learned agent gets to a known correct value, but whether the institution being measured was ever placed under real economic pressure in the first place. A property-rights rule is not being tested if no non-owner ever has reason to access a claimed resource; a trade price control is not being tested if every trade happens to be one-for-one already. Survival, sustainability, trade counts, property pressure, specialization, and institution-block counters exist to answer that prior question before any welfare number from this world is trusted. This is not a weaker test than the other four worlds run; it is a test of a condition the other four simply assume already holds, and Section 5 reports a case where that assumption failed and had to be corrected before the world's results meant anything.

\section{Platform}

\subsection{Interfaces}

The codebase is organized around three abstractions:

\begin{itemize}
  \item A \code{World} defines state transitions, observations, metrics, and
  rendering state.
  \item An \code{Agent} or mind maps observations to actions and updates from
  rewards and next observations.
  \item An \code{Institution} modifies rules, rewards, constraints, or
  information inside a world.
\end{itemize}

This separation is not cosmetic. It makes it possible to ask whether a result is
about a world, an institution, or a learner. It also creates a regression test:
if a new world requires rewriting the shared learning interface, then the
platform abstraction has failed.

\subsection{Validation Discipline}

The platform treats a result as reportable only after four checks. First, the
mechanics must be unit-tested: allocation, movement, payment, trade, collision,
or reward rules must match hand-computed cases. Second, where theory provides a
known answer, the implementation must recover or bracket that answer before
learned-policy output is interpreted. The current known-answer runner reports
19/19 passing checks across Pricing Arena N-firm static Nash references,
Auction House truthfulness and bid-shading cases, Public Goods free-rider and
cooperative brackets, Labor Market deferred-acceptance benchmark cases, and
Resource Island gather bounds. Third, the experiment must be replicated over
seeds, with per-seed and aggregate CSV outputs. Fourth, the institution must
actually activate. An institution that never faces the behavior it is meant to
govern produces no evidence about that institution, only about the world design.
Section 5.2 gives a concrete case. The second new validation layer is mechanism tracing. The trace runner produces
per-step rows and algebraic decomposition rows for all five worlds and all
thesis-facing minds. These traces do not replace full-run evidence; they show
the channel behind a claim and can expose contradictions between short traces
and full-run results. Section 4.3 gives a concrete case.

\section{Environment 1: Pricing Arena}

\subsection{World}

Pricing Arena is a repeated two-firm pricing market. Each firm chooses a price
from a discrete grid every round. Consumers divide between the two firms
according to a noisy logit demand model, so a firm's realized quantity depends
on both its own price and its rival's. Profit is price minus cost, times
realized quantity, and learning agents update from repeated play against each
other. Six mechanisms are implemented on top of this base game: no regulation,
a price cap, a tax on high prices, random audits, a direct anti-collusion
penalty, and demand shocks.

\subsection{Metrics}

The world reports average price, total profit, consumer surplus, welfare, price
dispersion, exploitability, and two collusion indexes. The first is a
price-normalized proxy used since early versions of the project. The second is
a profit-normalized index, used because it maps directly onto the collusion
measure common in the algorithmic-pricing literature:

\begin{equation}
  \mathrm{CI}_{profit}
  =
  \frac{\pi_{observed} - \pi_{Nash}}
       {\pi_{monopoly} - \pi_{Nash}}.
\end{equation}

Both are reported together, since the two can move in different directions, and
that divergence is itself informative about which channel a mechanism is
actually affecting.

\subsection{Known-Answer Validation}

Before any learned result is interpreted, the static Nash best response is
checked directly against the world's own demand and cost parameters, for
\(n=2\), \(3\), \(4\), and \(5\) firms. All four checks pass. This does not
validate anything about learned behavior; it validates that the reference point
learned behavior is measured against is not arbitrary.

\subsection{Baseline and Main Result}

Table~\ref{tab:pricing-baseline} shows the unregulated baseline for each
learner. Q-learning and the centralized critic sit at the higher-price,
higher-collusion end of the range; DQN, PPO, and independent-DQN sit lower,
with lower exploitability.

\begin{table}[h]
\centering
\begin{tabular}{lrrrrr}
\toprule
Mind & Avg.\ Price & Profit & Welfare & Collusion Index & Exploitability \\
\midrule
Q-learning & 5.487 & 365.3 & 810.4 & 0.541 & 95.7 \\
DQN & 3.611 & 236.3 & 832.8 & 0.202 & 20.3 \\
PPO & 3.681 & 260.0 & 839.3 & 0.215 & 14.2 \\
Independent-DQN & 4.078 & 269.9 & 827.9 & 0.285 & 26.6 \\
Centralized critic & 5.925 & 288.6 & 791.6 & 0.618 & 147.5 \\
\bottomrule
\end{tabular}
\caption{Pricing Arena, no-regulation baseline, by learner.}
\label{tab:pricing-baseline}
\end{table}

Table~\ref{tab:pricing-cap-full} shows the change under a price cap, relative
to that baseline, for every learner. The exploitability column is the clearest
result in this world: it falls for every learner, without exception. The profit
column is not: it falls for Q-learning and random play, in line with the
conventional expectation that a cap disciplines pricing behavior, but it rises
for DQN, independent-DQN, and the centralized critic.

\begin{table}[h]
\centering
\begin{tabular}{lrrrr}
\toprule
Mind & Price \(\Delta\) & Profit \(\Delta\) & Welfare \(\Delta\) & Exploitability \(\Delta\) \\
\midrule
Q-learning & \(-0.759\) & \(-22.916\) & \(+19.567\) & \(-55.366\) \\
Random & \(-1.185\) & \(-16.028\) & \(+34.779\) & \(-31.775\) \\
DQN & \(+0.388\) & \(+55.358\) & \(+6.294\) & \(-11.619\) \\
PPO & \(-0.042\) & \(-6.558\) & \(-1.198\) & \(-5.704\) \\
Independent-DQN & \(-0.288\) & \(+3.093\) & \(+12.183\) & \(-14.871\) \\
Centralized critic & \(-1.337\) & \(+10.347\) & \(+28.668\) & \(-98.274\) \\
\bottomrule
\end{tabular}
\caption{Price cap versus no regulation, change by learner, n=20 full run.}
\label{tab:pricing-cap-full}
\end{table}

The pattern is a guardrail that holds on the axis it was built for and does not
hold on an adjacent one. A price cap constrains the price level directly, and
every learner's exploitability falls as a result, since a capped price cannot
be pushed into the range a deviating opponent could exploit most severely. But
the cap does not constrain profit directly, only through price, and for three
of the five learners the resulting change in realized quantity and market
structure more than offsets the lower price. Q-learning behaves close to the
textbook intuition. DQN, independent-DQN, and the centralized critic do not.

\subsection{The Training-Budget Correction}

The DQN result above was not visible in the first version of this analysis. A
short, 1000-step mechanism trace, run to check the channel behind the headline
number, showed DQN profit falling under the cap for every seed tested, the
opposite of Table~\ref{tab:pricing-cap-full}. This contradiction was not
resolved by picking the number that fit the argument. It was resolved by
running the same paired-seed comparison at a range of training budgets, from
the short trace up to the full 40,000-step run the headline table is drawn
from.

\begin{table}[h]
\centering
\begin{tabular}{lrrrr}
\toprule
Source & Steps & Mean Profit \(\Delta\) & Positive Seed Share & Mean Quantity \(\Delta\) \\
\midrule
fresh audit & 1{,}000 & \(-11.711\) & 0.10 & 3.220 \\
fresh audit & 5{,}000 & \(+1.730\) & 0.50 & 4.102 \\
fresh audit & 10{,}000 & \(+35.448\) & 0.80 & 3.540 \\
full table & 40{,}000 & \(+55.358\) & 0.85 & 0.226 \\
\bottomrule
\end{tabular}
\caption{DQN price-cap profit delta by training budget. Values are paired seed
deltas, price cap minus no regulation.}
\label{tab:pricing-budget-ladder}
\end{table}

The sign flips between 1,000 and 5,000 steps and is positive with a clear
majority of seeds by 10,000. At 40,000 steps, the profit effect is
statistically positive (\(\Delta\pi = 55.358\), 95\% CI
\([4.859,\ 105.857]\), positive in 85\% of seeds), but the mean quantity
change at that budget is small, unlike the 5,000 and 10,000-step rows where
quantity moves substantially. This rules out the simplest explanation, that
DQN is exploiting a fixed quantity-margin loophole from the start; the effect
is not present early and does not have a single stable channel behind it at
every budget. The safer claim is that the price cap becomes a focal constraint
some DQN runs learn to use to their advantage only after enough training, not a
channel visible or exploitable from the beginning. A short trace is not
sufficient evidence either way for a result like this one; the training-budget
ladder is what makes the finding reportable.

\subsection{The Other Four Mechanisms}

Price cap is the only mechanism in this world with a result this clean, and
that asymmetry is itself worth reporting rather than hiding. The remaining four
are correctly implemented and pass the same unit and parity checks as price
cap, but three of the four are weak, inert, or inconsistent enough across
learners that no single clear claim follows from them at the currently tested
parameters.

The high-price tax subtracts a quantity-weighted penalty from reward above a
price threshold. Its effect is small in every direction and every learner:
Q-learning's exploitability falls modestly, DQN is essentially unaffected, and
the centralized critic moves in the wrong direction on price and welfare. At
the tested tax rate, the penalty is too small, or too easy to avoid, to reliably
change pricing behavior.

\begin{table}[h]
\centering
\begin{tabular}{lrrr}
\toprule
Mind & Price \(\Delta\) & Profit \(\Delta\) & Exploitability \(\Delta\) \\
\midrule
Q-learning & \(-0.078\) & \(+2.2\) & \(-7.6\) \\
DQN & \(+0.006\) & \(+0.6\) & \(+1.2\) \\
PPO & \(-0.023\) & \(-3.6\) & \(-5.2\) \\
Independent-DQN & \(+0.001\) & \(+0.2\) & \(+2.5\) \\
Centralized critic & \(+0.088\) & \(+4.2\) & \(-1.2\) \\
\bottomrule
\end{tabular}
\caption{High-price tax versus no regulation, change by learner.}
\label{tab:pricing-tax}
\end{table}

Random audits apply a stochastic penalty when average price is high. The result
here is worse than merely weak: for Q-learning, exploitability rises under
audits rather than falling, and for the centralized critic, profit rises
sharply while exploitability falls only partially. An enforcement mechanism
that increases exploitability for one learner while partially working for
another is not a guardrail; at these parameters it is closer to noise with an
inconsistent sign.

\begin{table}[h]
\centering
\begin{tabular}{lrrr}
\toprule
Mind & Price \(\Delta\) & Profit \(\Delta\) & Exploitability \(\Delta\) \\
\midrule
Q-learning & \(+0.114\) & \(+11.1\) & \(+32.0\) \\
DQN & \(+0.035\) & \(+4.3\) & \(-0.5\) \\
PPO & \(-0.006\) & \(-1.0\) & \(-2.7\) \\
Independent-DQN & \(+0.002\) & \(+0.2\) & \(+1.7\) \\
Centralized critic & \(+0.125\) & \(+49.1\) & \(-21.7\) \\
\bottomrule
\end{tabular}
\caption{Random audit versus no regulation, change by learner.}
\label{tab:pricing-audit}
\end{table}

The direct anti-collusion penalty punishes prices that are both high and close
together, the most explicit mechanism in the suite and, on paper, the one
closest to actual antitrust enforcement logic. It still mostly fails to bind:
four of the five learners show negligible or worsened exploitability, and only
PPO shows a mild, coherent improvement across price, profit, and exploitability
together. A rule aimed directly at the behavior of interest is not
automatically more effective than the blunter price cap.

\begin{table}[h]
\centering
\begin{tabular}{lrrr}
\toprule
Mind & Price \(\Delta\) & Profit \(\Delta\) & Exploitability \(\Delta\) \\
\midrule
Q-learning & \(-0.048\) & \(-2.1\) & \(+31.9\) \\
DQN & \(+0.005\) & \(+0.5\) & \(0.0\) \\
PPO & \(-0.067\) & \(-7.4\) & \(-5.5\) \\
Independent-DQN & \(+0.001\) & \(+0.2\) & \(0.0\) \\
Centralized critic & \(-0.013\) & \(+0.1\) & \(+0.5\) \\
\bottomrule
\end{tabular}
\caption{Anti-collusion penalty versus no regulation, change by learner.}
\label{tab:pricing-anticollusion}
\end{table}

Demand shocks perturb market size with a lognormal multiplier each round,
testing robustness to environmental volatility rather than acting as a policy
lever. Results are small and mixed in direction across learners, with the
centralized critic the only one meaningfully affected, and negatively. This
mechanism is best read as a stress test the platform passes mechanically, not
as a candidate institutional finding.

\begin{table}[h]
\centering
\begin{tabular}{lrrrr}
\toprule
Mind & Price \(\Delta\) & Profit \(\Delta\) & Welfare \(\Delta\) & Exploitability \(\Delta\) \\
\midrule
Q-learning & \(+0.014\) & \(+4.0\) & \(+0.3\) & \(+11.4\) \\
DQN & \(-0.002\) & \(+1.3\) & \(+1.8\) & \(+1.3\) \\
PPO & \(-0.010\) & \(-0.3\) & \(+1.4\) & \(-6.0\) \\
Independent-DQN & \(+0.001\) & \(+0.6\) & \(+1.2\) & \(+1.8\) \\
Centralized critic & \(-0.050\) & \(-9.2\) & \(-12.0\) & \(+2.7\) \\
\bottomrule
\end{tabular}
\caption{Demand shock versus no regulation, change by learner.}
\label{tab:pricing-demandshock}
\end{table}

None of the four should be read as a code defect. All four pass the same
mechanics and parity tests as price cap. The difference is economic activation,
not implementation: at the parameters currently tested, only price cap creates
pressure strong enough to produce a consistent, cross-learner, interpretable
signal. Table~\ref{tab:pricing-mechanism-status} summarizes the distinction.

\begin{table}[h]
\centering
\begin{tabular}{lll}
\toprule
Mechanism & Implementation & Interpretability \\
\midrule
Price cap & Passes all checks & Strong, consistent, cross-learner \\
High-price tax & Passes all checks & Weak at tested rate \\
Random audit & Passes all checks & Mixed, sign-inconsistent \\
Anti-collusion penalty & Passes all checks & Mostly inert, brittle \\
Demand shock & Passes all checks & Robustness check, not a policy result \\
\bottomrule
\end{tabular}
\caption{Mechanism status, implementation versus interpretability, Pricing Arena.}
\label{tab:pricing-mechanism-status}
\end{table}

\subsection{Summary}

Pricing Arena's clearest result is a guardrail that succeeds on the axis it
directly targets and fails to hold as a guarantee on an adjacent one: price caps
lower exploitability for every learner tested, but their effect on profit and
collusion is learner-dependent, with the effect for DQN emerging only after
sufficient training rather than being present from the start. The remaining
four mechanisms are correctly implemented but, at current parameters, mostly
too weak or inconsistent to support a claim beyond their own numbers. Whether
the price-cap result and the weakness of the other four persist as the number
of firms increases beyond two is addressed in the cross-world synthesis
section.

\section{Environment 2: Resource Island}

\subsection{World}

Resource Island is a bounded-grid gather/trade economy. Agents move, gather food
and wood, maintain energy, accumulate inventories, and can die from starvation.
The world contains property rights, redistribution, trade price controls, and
reputation institutions. It was built to test whether the platform abstractions
work outside a two-firm pricing game.

\subsection{What v0 Taught}

The first full Resource Island runs were mechanically valid but economically
under-activated. In the initial full run, successful trade was sparse and some
institutions did not bind. Follow-up replay showed that property rights almost
never faced behavioral pressure: even when claims existed, non-owner
opportunities to gather from claimed resource cells were extremely rare.
Similarly, trade price controls could not bind because the original trade
protocol only allowed one-for-one exchange.

This is an important negative result for the platform: a world can be runnable,
tested, and statistically replicated while still failing to exercise the
institution it is supposed to study.

\subsection{v1 Hardening}

Resource Island v1 adds four changes, motivated by common-pool-governance and
public-goods work on monitoring, local fit, incentives, and repeated
interaction \citep{ostrom1990,fehrgachter2000,isaac1994}:

\begin{enumerate}
  \item contested resource layouts, so multiple agents have reason to access the
  same cells;
  \item unequal trade offers, so price controls can actually bind;
  \item specialization pressure through heterogeneous resource preferences and
  starting inventories;
  \item activation diagnostics for trade attempts, blocked trades, property
  opportunities, property violations, and institution blocks.
\end{enumerate}

The v1 validation run verified that the economic channels were active, and the
subsequent n=20 full run preserves that activation. Table~\ref{tab:resource-v1}
shows the medium validation values used as the pre-full-run gate.

\begin{table}[h]
\centering
\begin{tabular}{lrrrrr}
\toprule
Institution & Welfare & Survival & Trades & Trade Blocks & Property Opps. \\
\midrule
none & 1.095 & 0.933 & 6.034 & 0.000 & 0.000 \\
property rights & 1.078 & 0.923 & 5.790 & 0.000 & 17.734 \\
trade price controls & 0.875 & 0.898 & 0.000 & 5.779 & 0.000 \\
reputation & 1.313 & 0.927 & 6.867 & 0.000 & 0.000 \\
\bottomrule
\end{tabular}
\caption{Resource Island v1 validation results.}
\label{tab:resource-v1}
\end{table}

The interpretation is straightforward. Price controls now bind because the
world permits unequal trades. Reputation increases successful trade and welfare
in the validation setting. Property-right opportunities are present, although
the welfare effect is small in this medium run. The full n=20 v1 run confirms
the activation pattern: successful trade appears under no institution, price
controls bind and suppress trade, property-right opportunities are measured,
and reputation produces the strongest welfare gains in this pressure setup.
The strict-local ablation confirms that this remains partly a spatial-friction
story: reducing trade from all-island matching to radius-one local matching cuts
baseline trade from 5.325 to 1.900 and reputation trade from 6.667 to 1.927,
while property opportunities and price-control blocks remain positive.

\subsection{Cross-Mind Resource Island}

The v0 Resource Island Phase 3 full table initially gave a different pattern
from Pricing Arena: neural/MARL minds improved survival and welfare but did not
learn successful trade. After the v1 pressure changes, that is no longer the
main Resource Island conclusion. Table~\ref{tab:resource-v1-minds} shows the
no-institution row from the v1 full learner-suite run.

\begin{table}[h]
\centering
\begin{tabular}{lrr}
\toprule
Mind & Welfare & Trades \\
\midrule
Q-learning & 1.107 & 5.325 \\
DQN & 1.061 & 0.640 \\  
PPO & 1.254 & 8.892 \\
Independent DQN & 1.063 & 1.730 \\
Centralized critic & 1.248 & 8.697 \\
\bottomrule
\end{tabular}
\caption{Resource Island v1 full learner-suite run, no-institution condition.}
\label{tab:resource-v1-minds}
\end{table}

The updated result is more useful. PPO and centralized-critic discover the
trade channel under pressure, while DQN and fixed independent-DQN remain more
conservative and trade much less. Under trade price controls, successful trade
falls to zero across the learner suite because unequal exchange is blocked. Resource
Island therefore contributes a different kind of learner-architecture effect from Pricing
Arena: stronger learners do not simply exploit or avoid an institution
monotonically; their ability to activate the intended market channel depends on
the learning rule and on whether the environment creates enough pressure for
exchange.

\section{Environment 3: Auction House}

\subsection{World}

Auction House is a single-item sealed-bid auction world with independent private
values, discrete bids, deterministic tie-breaking, first-price and second-price
payment rules, and optional reserve prices \citep{vickrey1961,myerson1981,milgrom2004}. Each bidder observes its own
valuation bin and chooses a bid. The world reports seller revenue, bidder
surplus, total welfare, allocative efficiency, regret, overbidding, underbidding,
and distance from benchmark bid functions.

\subsection{Benchmarks}

The second-price benchmark is truthful bidding \citep{vickrey1961}. The first-price benchmark is
bid shading under symmetric independent private values, with discrete-grid
regret checks used where the continuous benchmark is too coarse
\citep{milgrom2004,dutting2019,banchio2022}. Reserve-price
variants evaluate the revenue-efficiency tradeoff against the no-reserve
benchmark \citep{vickrey1961,myerson1981,milgrom2004}.

\subsection{Validation}

The Auction House validation run is economically coherent but not yet final
evidence. Second-price learning moves toward the truthful benchmark, first-price
learning shows underbidding and bid shading, and reserve prices increase
revenue while changing allocative efficiency. Table~\ref{tab:auction-validation}
summarizes the current validation run.

\begin{table}[h]
\centering
\begin{tabular}{lrrrr}
\toprule
Scenario & Revenue & Welfare & Efficiency & Underbid Rate \\
\midrule
second price & 3.373 & 6.216 & 0.747 & 0.381 \\
first price & 3.481 & 6.331 & 0.768 & 0.774 \\
second price + reserve & 4.031 & 6.098 & 0.814 & 0.366 \\
\bottomrule
\end{tabular}
\caption{Auction House validation results.}
\label{tab:auction-validation}
\end{table}

\subsection{Cross-Mind Auction Diagnostics}

The full P.6 auction learner-suite run confirms that the auction world should be read
against benchmark deviations, not raw reward alone. In second-price auctions,
PPO has the lowest ex-post regret among the tested learning minds
(\(0.187\)) and higher allocative efficiency than Q-learning (\(0.762\)
versus \(0.734\)). DQN and fixed independent-DQN earn higher seller revenue but
with lower allocative efficiency. The centralized-critic scaffold performs
poorly in this asymmetric bidding setting, with low revenue and high regret.

First-price auctions show the expected underbidding/shading channel across the
non-centralized learners \citep{milgrom2004}. Q-learning, DQN, PPO, and fixed independent-DQN all
underbid frequently, with first-price underbid rates between \(0.716\) and
\(0.814\) \citep{milgrom2004}. This makes Auction House theory-anchored: the economic question is
how each learner deviates from truthful second-price bidding or shaded
first-price bidding, not whether the environment merely produces positive
payoffs \citep{vickrey1961,milgrom2004,dutting2019}.

\section{Environment 4: Public Goods}

\subsection{World}

Public Goods is a repeated common-pool resource environment, following the
standard tension between private extraction and shared surplus
\citep{samuelson1954,hardin1968}. Agents choose
between contributing to the pool, extracting from it, or remaining inactive.
Extraction pays private reward but depletes the shared stock; contribution costs
the contributor but helps sustain the pool. When requested extraction exceeds
available stock, extraction is rationed proportionally. Deterministic
regeneration makes the sustainability channel easy to audit.

The implemented institutions are no intervention, a penalty schedule,
contribution matching, reputation rewards, information restriction, and a
generic tax-and-redistribution schedule. The main metrics are welfare,
sustainability, total contribution, extraction relative to regeneration,
inequality, collapse diagnostics, and tax revenue where applicable.

\subsection{Q-Learning Full Run}

The n=20 tabular Q-learning full run confirms the intended commons pressure. In
the baseline, agents mostly extract and contribute little: sustainability is
0.0892 and total contribution is 0.1418. Contribution matching improves both
welfare and sustainability in this setting, reaching welfare 2.0808 and
sustainability 0.1048. Reputation is a strong reward-shaping intervention, with
welfare 9.5750, while the tax schedule mainly changes redistribution/accounting
at the tested rates.

The effect validator classifies the tax schedule as mostly reward/accounting
under this configuration, while penalty, contribution matching, reputation, and
information restriction change state metrics relative to baseline. This is a
diagnostic classification, not yet a final welfare ranking.

\subsection{Cross-Mind Status}

The full Public Goods learner suite shows that stronger function approximation
or centralized training does not automatically solve the commons problem. Under
no institution, Q-learning reaches welfare \(1.815\),
sustainability \(0.089\), and contribution \(0.142\). DQN and fixed
independent-DQN are close but contribute less; PPO and centralized-critic
contribute almost nothing and sit near the lower sustainability regime.

Contribution matching is the clearest state-changing intervention in the full
learner-suite run. It improves Q-learning sustainability from \(0.089\) to \(0.105\) and
fixed independent-DQN sustainability from \(0.085\) to \(0.121\), with a larger
increase in contribution for independent-DQN. It does not rescue PPO in this
configuration. The economic reading is therefore conditional: the institution
can improve the state variables when the learner explores contribution, but a
more capable policy class can still converge to low-contribution behavior.

\section{Environment 5: Labor Market}

\subsection{World}

Labor Market is a two-sided matching environment. Workers are learning agents;
employers have fixed preferences. Each worker reports a top choice, the world
constructs reported preferences, and matching is resolved using deferred
acceptance \citep{galeshapley1962,rothsotomayor1990}. The world reports match rate, total welfare, truthful-report rate,
stability, blocking-pair diagnostics, and manipulation-gain diagnostics.

This world is intentionally asymmetric. Unlike Pricing Arena, Resource Island,
Public Goods, and Auction House, not every economic actor is a learning agent.
That asymmetry is the main interface test for the learning-architecture suite:
DQN, PPO, independent-DQN, and centralized-critic operate only over the worker
side while employers remain fixed preference holders.

\subsection{Q-Learning Full Run}

The n=20 tabular Q-learning full run is substantially more stable than the
short smoke run. Match rate remains \code{1.0}, while
\code{stability\_mean=0.9508}, \code{truthful\_report\_rate\_mean=0.7859}, and
\code{total\_welfare\_mean=3.6393}. The report-top action space creates
truthfulness variation, but worker-proposing deferred acceptance still produces
mostly stable matchings under the current random-preference setup.

Fixed benchmark cases clarify the interpretation. Under worker-proposing
deferred acceptance, workers should not have profitable report deviations in
the canonical strategy-proof setting \citep{galeshapley1962,rothsotomayor1990,roth1984}. Future manipulation tests therefore need
to target the non-proposing side, change the mechanism, or add information and
commitment frictions rather than treating worker misreports under standard
deferred acceptance as the main expected failure mode.

\subsection{Cross-Mind Status}

The full Labor Market learner-suite run keeps match rate at \(1.0\) for every mind, so the
relevant margins are stability and truthful-report behavior. DQN slightly
improves stability relative to Q-learning (\(0.975\) versus \(0.951\)) while
reporting truthfully less often. PPO and fixed independent-DQN preserve high
stability but also lower truthfulness. The centralized-critic scaffold has the
weakest matching behavior in this asymmetric setup, with stability \(0.749\)
and truthful-report rate \(0.433\).

Economically, this is consistent with the benchmark warning: worker-proposing
deferred acceptance is strategy-proof for workers in the canonical setting, so
lower truthful-report rates are not automatically evidence of profitable
manipulation \citep{galeshapley1962,rothsotomayor1990,roth1984}. The result is better framed as a learned-reporting diagnostic:
the mechanism keeps matches mostly stable for most learners, while the
centralized-critic architecture is brittle in this asymmetric environment.

\section{Cross-World Synthesis}

The emerging thesis is not that one institution always works, or that the
agent classes form a universal intelligence ranking. The stronger claim is more
specific: institutional robustness must be evaluated jointly with the learning
architecture, the training horizon, and activation diagnostics. Tabular value
learning, function approximation, on-policy policy gradients, decorrelated
independent learners, and centralized critics violate different assumptions of
the classical benchmark story. The comparisons are therefore best read as a
learner-architecture stress test, not as a scalar capability ranking.

The Pricing Arena result shows architecture-sensitive institutional effects:
price caps reduce exploitability broadly, but DQN's profit response flips sign
with training budget and is positive in the full paired audit. Resource Island
shows the complementary failure mode: if the world design does not put an
institution under pressure, then the absence of an effect is not a mechanism
result. Auction House adds theory-anchored bidding benchmarks
\citep{vickrey1961,myerson1981}. Public Goods adds a sustainability and
commons-collapse channel. Labor Market adds a matching-stability and
truthfulness channel.

The five-world synthesis is now built from full n=20 learner-suite outputs for Pricing
Arena, Resource Island v1, Auction House, Public Goods, and Labor Market. The
protocol-comparability report passes for all five worlds, and the publication
inference table includes all worlds under the same confirmatory bookkeeping.
The claim-audit suite adds a second pass over those outputs: paired seed
effects, benchmark gaps, metric-conflict warnings, activation confirmations,
and trace/full-run sign agreement where paired traces exist. In its current
run, the audit produces 1,025 paired or benchmark rows and 31 metric-conflict
or activation rows. This audit layer prevents the cross-world synthesis from
becoming a table of averages without claim boundaries.

The main synthesis claim is qualitative but concrete. Learning architecture
changes institutional outcomes in every world, but not through one universal
direction. In Pricing Arena, DQN's price-cap profit increase is a
long-training, seed-heterogeneous effect rather than a short-run channel. In
Resource Island v1, PPO and centralized-critic activate trade strongly while
DQN-family learners remain more conservative \citep{mnih2015,schulman2017}. In Auction House, the relevant
question is distance from truthful, shaded, or reserve-price benchmarks rather
than raw reward \citep{vickrey1961,myerson1981,milgrom2004}. In Public Goods, the relevant state variables are
contribution, extraction, and sustainability, while some tax effects are
accounting effects rather than state changes. In Labor Market, deferred
acceptance produces high match rates and mostly stable outcomes, so future
manipulation claims must target the right side of the market or a different
matching institution.

\section{Limitations}

Several limitations are deliberate. The environments are small and stylized.
They are designed for auditability rather than realism. The deep-RL agents are
structured-observation MLP agents rather than large-scale policies. The
Resource Island trade protocol is still simple, and the current v1 pressure
configuration is an activation test rather than a realistic market. Auction
House currently has benchmarked sealed-bid full runs and smoke-tested clock and
information variants. Public Goods is still a stylized common-pool model rather
than an empirical commons. Labor Market currently uses worker-side learning
under deferred acceptance, so it tests one asymmetric matching channel rather
than matching-market design generally.

The most important limitation is interpretive. A clean CSV does not imply a
clean economic result. Every institution must be checked for activation. Every
learner-architecture comparison must be checked for implementation artifacts. The
independent-DQN alias issue is a concrete example: the table was statistically
valid as an output file, but not valid as evidence for a distinct mind until the
implementation was corrected and rerun.

The Pricing Arena price-cap audit adds a second interpretive warning. A
mechanism trace can reveal the local channel but cannot replace a replicated
training-budget audit. The 1000-step trace correctly showed early profit losses
under the cap; the 40,000-step table correctly showed a positive DQN
cap-minus-none profit delta. The result is reportable only because the
contradiction was resolved through paired seed deltas and a budget ladder.

\section{Reproducibility}

The repository records validation status in \code{ROADMAP\_STATUS.md}. Core
outputs are stored under \code{outputs/}. Major runners include:

\begin{itemize}
  \item \code{run\_multiseed.py} for Pricing Arena replicated runs;
  \item \code{run\_exploitability.py} for frozen-policy adversarial evaluation;
  \item \code{run\_phase3\_full.py} and \code{validate\_phase3\_full.py} for
  full Phase 3 Pricing Arena checks;
  \item \code{run\_known\_answer\_sanity\_checks.py} for theorem/bracketing
  checks before learned-policy interpretation;
  \item \code{run\_mechanism\_traces.py} for per-step mechanism traces and
  decomposition tables across worlds;
  \item \code{run\_pricing\_quantity\_margin\_audit.py} for the Pricing Arena
  price-cap budget-ladder contradiction audit;
  \item \code{run\_claim\_audit\_suite.py} for cross-world paired-effect,
  benchmark-gap, metric-conflict, activation, and trace-alignment audits;
  \item \code{run\_resource\_island\_smoke.py} for Resource Island smoke,
  validation, and full runs;
  \item \code{run\_auction\_house\_smoke.py} for Auction House runs;
  \item \code{run\_public\_goods\_smoke.py} for Public Goods runs;
  \item \code{run\_labor\_market\_smoke.py} for Labor Market runs;
  \item \code{build\_world\_mind\_comparison.py} and
  \code{build\_cross\_world\_synthesis.py} for ladder comparisons and
  cross-world synthesis.
\end{itemize}

The test suite currently includes unit, integration, and output-validation
coverage for the core abstractions, Pricing Arena, Resource Island, Auction
House, Public Goods, Labor Market, and deep-RL/MARL minds.

\subsection{Primary Output Artifacts}

The main numerical claims in the paper are backed by the following generated
artifacts:

\begin{itemize}
  \item \path{outputs/known_answer_sanity_checks_current/summary.csv};
  \item \path{outputs/mechanism_traces_local_full/mechanism_decomposition.csv};
  \item \path{outputs/phase3_full/mind_comparison.csv};
  \item \path{outputs/pricing_quantity_margin_audit/dqn_budget_ladder.csv};
  \item \path{outputs/pricing_quantity_margin_audit/existing_full_seed_deltas.csv};
  \item \path{outputs/resource_island_v1_full/summary_aggregate.csv};
  \item \path{outputs/resource_island_v1_strict_radius_full/summary_aggregate.csv};
  \item \path{outputs/resource_island_v1_phase3_full/mind_comparison.csv};
  \item \path{outputs/auction_house_phase3_full/mind_comparison.csv};
  \item \path{outputs/public_goods_phase3_full/mind_comparison.csv};
  \item \path{outputs/labor_market_phase3_full/mind_comparison.csv};
  \item \path{outputs/cross_world_synthesis/synthesis_table.csv};
  \item \path{outputs/claim_audit_suite/claim_audit_report.md}.
\end{itemize}

\section{Conclusion}

This draft reports a first version of an artificial-economies platform for
testing institutional robustness under learning agents. The mature Pricing
Arena result already shows why the question matters: exploitability and
collusion can move differently under the same institution, and a headline
effect can depend on training budget rather than appearing in short traces.
Resource Island, Auction House, Public Goods, and Labor Market extend the
project beyond pricing games. Their role is not merely to add more checkboxes,
but to test whether the patterns found in Pricing Arena survive when the
economic structure changes and when the relevant institutional channel is
actually activated.

\bibliographystyle{plainnat}
\bibliography{references}

\end{document}
